\section{評価と考察}

\subsection{評価課題と実験条件}
本稿では，提案手法単体の関数レベルの判定能力を評価する．実際のTPL検出器が生成した候補集合や，ライブラリ単位の検出精度向上は評価対象としない．次のResearch Question（RQ）に沿って評価する．
\begin{description}
  \item[RQ1] 等価ペアと非等価ペアを，それぞれどの範囲で安全にEQまたはNOT\_EQと判定できるか．
  \item[RQ2] 制御構造，状態変化およびビルド差と判定結果には，どのような関連があるか．
  \item[RQ3] UNKNOWNとTIMEOUTは，どの証拠不足または計算上の限界から生じるか．
  \item[RQ4] 関数レベルの結果は，指定バージョンとの厳密な対応候補を選別する後段検証へどのように利用できるか．
\end{description}

\paragraph{等価ペアの評価}
等価セットでは，同一意味を持つ基準ソースを16条件でビルドし，生成されたバイナリ間のEQ判定能力を評価する．EQBENCH上の基準ソース名は\texttt{new.c}または\texttt{newV.c}である．ビルド条件は，命令セット2種（ARM64，MIPS64），コンパイラ2種（Clang，GCC），最適化4段階（O0--O3）の直積であり，1ソース当たり16バイナリを生成する．1関数について異なる2バイナリの全組合せを比較するため，比較数は$16 \times 15 / 2 = 120$ペアである．現在のEqログは，この120ペアの結果を原関数ごとに記録したものである．

同一ソース由来であることだけを等価性の根拠とはせず，未定義動作，処理系定義動作，整数・ポインタ幅および符号付きオーバーフローへの依存を確認し，比較する入力領域と観測結果の幅を明示する．

\paragraph{非等価ペアの評価}
非等価セットでは，意味を変更する前後のソースを新側・旧側として比較し，NOT\_EQ判定能力を評価する．EQBENCH上では，\texttt{new.c}と\texttt{old.c}，または\texttt{newV.c}と\texttt{oldV.c}の組を用いる．定数・演算子，条件境界，分岐出力，ループ更新，メモリ更新または呼出し効果など，導入された意味差の種類も記録する．Neq実験は現在検証中であり，組合せ方法と最終比較数が確定するまでは，NOT\_EQ判定能力の結果を確定値として記載しない．

% TODO(neq-design): 新側16種と旧側16種の全組合せを比較するか，同一ビルド条件だけを対応させるか，new--oldとnewV--oldVを別集計するか，および1関数当たりの最終比較数を，実験完了後に確定する．

\paragraph{共通の評価単位}
EqとNeqは母数と組合せ規則が異なるため別集計とし，全実験を合算した指標は補助情報に限定する．各評価系では，ペア加重集計とベンチマークケース平均を併記する．ペア加重集計は全比較ペアを数える．ベンチマークケース平均では，Eqは基準ソース中の1関数を1ケース，Neqは意味変更前後の関数対を1ケースとし，各ケースの判定率を計算してから平均する．ビルド差の影響は同一ケース内で比較し，関数または変更ケース自体の難しさとコンパイラ・命令セット差を分離する．ERROR，未完了および成果物欠損は実験健全性として別に集計し，対象から黙って除外しない．

\paragraph{実行条件}
EqとNeqに共通して，解析プログラムの版，実行環境，並列数，ペア当たりの制限時間およびソルバ設定を固定する．

% TODO(experiment): EQBENCHの版と等価性の定義，未定義・処理系依存動作の除外条件，入力領域・出力幅の対応，対象関数数，解析プログラムの版，並列数，制限時間，実行環境および再現手順を，最終成果物の識別情報とともに固定する．

\subsection{判定能力（RQ1）}
4値判定を単一のaccuracyへまとめず，等価ground truthと非等価ground truthに分けて示す．等価セットでは，EQ，NOT\_EQ，UNKNOWN，TIMEOUTの件数と割合を示し，EQ率と判定網羅率を主要指標とする．等価ペアをNOT\_EQとした結果は危険な誤判定として独立に報告する．非等価セットでも4値の件数と割合を別表で示し，NOT\_EQ率と判定網羅率を主要指標とする．非等価ペアをEQとした結果は危険な誤判定である．判定網羅率は，各セットにおける「（EQ件数＋NOT\_EQ件数）／対象ペア数」とする．

本評価では，解析が生成した入力割当てを対象バイナリへ与える独立実行は行わない．したがって，解析内部で差の安全性を確認した結果を「反例確認率」として重ねて数えず，NOT\_EQに対して入力割当てと観測差の根拠を追跡可能な成果物として保存できた割合を，反例成果物生成率として補助的に示す．これはNOT\_EQの正しさを別経路で検証する指標ではなく，判定根拠の監査可能性を表す．

UNKNOWNおよびTIMEOUTは，誤ったEQまたはNOT\_EQと同じ失敗として扱わない．差が存在する非等価ペアをUNKNOWNとした場合は，非等価性を否定したのではなく，安全な反例として確定するための証拠が不足したものとして分類する．EqとNeqを合算した判定網羅率は，ペア数の多い評価系に支配されるため主要な結論には用いない．危険な誤判定数と判定網羅率を並べることで，保守性と適用範囲の両方を評価する．

% TODO(result-eq): 完了したEq成果物から，ペア加重集計とベンチマークケース平均による4値判定，EQ率，判定網羅率および危険なNOT\_EQ件数を作成する．
% TODO(result-neq): Neq実験完了後に，ペア加重集計とベンチマークケース平均による4値判定，NOT\_EQ率，判定網羅率，反例成果物生成率および危険なEQ件数を作成する．

\subsection{判定を分ける特徴と保留原因（RQ2・RQ3）}
制御構造は，直線的な処理，条件分岐，複数の終了経路，ループおよび再帰に分ける．状態変化は，戻り値のみ，ポインタ出力，メモリ更新および外部関数呼出しに分ける．ビルド差は，コンパイラ，最適化，命令セットおよび値の幅に分け，分岐数，論理式の規模および反復関係の有無も記録する．これらと判定結果のクロス集計はEqとNeqを分けて示す．Eq側のUNKNOWNは等価性を証明するための証拠不足，Neq側のUNKNOWNは観測差を安全な反例として確定するための証拠不足を表すためである．EqとNeqの合算値は補助情報に限定する．クロス集計からは「関連が観測された」と述べるにとどめ，原因は代表例の解析と同一ベンチマークケース内の比較によって説明する．

UNKNOWNについては，最終判定がUNKNOWNとなった比較数を「UNKNOWNペア数」，各ペアに記録された証拠不足の延べ数を「原因出現数」，同一ペア内で同じ分類に属する原因の重複を除いたペア数を「根本原因別ペア数」として区別する．原因は第3節の5分類へ対応づける．1ペアが複数分類の原因を持つことを許す複数回答形式であるため，原因ごとの件数をUNKNOWNペア数で割った割合の合計は100\%になるとは限らない．実装内部の詳細な識別名の延べ数は主結果とせず，5分類への対応を再現するための補助情報とする．

終了経路が増える場合は経路条件の網羅，ポインタ出力や呼出しを含む場合は副作用の把握，異なる命令セット間では値幅と呼出し規約の対応，ループ・再帰では反復関係の証明が判定の境界となり得る．必要な振る舞いを論理式へ変換できない場合はUNKNOWN，式を構成できても制限時間内に解けない場合はTIMEOUTである．結果確定後は，これらの一般論を観測結果と代表例へ統合し，重複する説明を削除する．

% TODO(analysis): EqとNeqを分け，各特徴について，ベンチマークケース，ビルド条件，観測結果，不足した証拠または計算量を同じ順序で示す．UNKNOWNはペア数，原因出現数，根本原因別ペア数を区別し，TIMEOUTは設定時間と処理時間の分布を併記する．

\subsection{後段検証への適用可能性と妥当性への脅威（RQ4）}
関数レベル評価で直接示せるのは，Eqセットでどの範囲の等価ペアを安全にEQとできるか，およびNeqセットでどの特徴を持つ非等価ペアを安全にNOT\_EQとできるかである．NOT\_EQは，当該バージョンの関数とは観測結果が異なることを示すため，正確なバージョン同定を必要とするSBOM作成に有用な負の証拠となる．一方，同じTPLの別バージョンに由来する可能性や，対象バイナリにTPL自体が存在する可能性は否定しない．EQも観測モデルの範囲で関数対応を支持する結果であり，単独でライブラリの同一性を証明するものではない．

したがって，本評価は，TPL検出システム全体のprecisionまたはrecallが向上したことを裏付けない．ライブラリ判定の偽陽性低減を定量化するには，実際のTPL検出器が生成した候補集合へ接続し，関数候補の除外前後で，バージョンを含むライブラリ単位のground truthと比較する統合評価が必要である．

妥当性への脅威として，EQBENCHが実際のTPL候補に現れる関数分布を代表するとは限らないこと，同一ソースでも未定義・処理系依存動作により意味が変わり得ること，観測モデル外の副作用があることを挙げる．現在の主要課題は，終了経路の網羅，呼出し先・メモリ効果の要約，一般的なループ・再帰の反復関係，異なる値幅の対応，反例の安全性確認および論理式の計算量である．これらはUNKNOWNまたはTIMEOUTの原因と対応づけ，関数レベル評価の適用限界として示す．
